% Source: https://artofproblemsolving.com/community/c3h458403_1997_mc_chapter_problems
% Stable identity: aops-topic-458403-q11
Let the square of tin have side $s$. Cutting a $1$-inch square from each corner and folding up the four flaps leaves a base that is a square of side $s-2$ (both dimensions lose $1$ inch at each end), and the height of the tray is $1$ inch.

\begin{asy}
usepackage("amsmath");
size(5cm);
defaultpen(fontsize(10pt));
real s=8;
fill((1,1)--(7,1)--(7,7)--(1,7)--cycle,paleyellow);
draw((0,0)--(s,0)--(s,s)--(0,s)--cycle,linewidth(1));
for(pair c: new pair[]{(0,0),(7,0),(7,7),(0,7)}) fill(shift(c)*unitsquare,gray(0.75));
for(pair c: new pair[]{(0,0),(7,0),(7,7),(0,7)}) draw(shift(c)*unitsquare,linewidth(0.8));
draw((1,1)--(7,1)--(7,7)--(1,7)--cycle,dashed);
draw((1,0)--(1,1),linewidth(0.8));draw((7,0)--(7,1),linewidth(0.8));
label("$1$",(0.5,0),S);label("$1$",(0,0.5),W);
label("$s-2=6$",(4,4));
draw((0,-1.1)--(s,-1.1),Arrows(TeXHead));label("$s=8$",(4,-1.1),S);
label("\mbox{}",(min(currentpicture,user=true)+max(currentpicture,user=true))/2,black); // keeps the text after the figure black when LaTeX typesets the labels inline
\end{asy}

So the base is not just any rectangle with area $36$: it is a \emph{square}, because the original piece is a square and the same amount is removed on every side. Hence
\[
(s-2)^2\cdot1=36\qquad\Longrightarrow\qquad s-2=6\qquad\Longrightarrow\qquad s=8,
\]
and the original square has area $8^2=\boxed{64}$ square inches. (The dashed lines in the figure are the folds.)
